\chapter{Preface}

This book teaches the Go programming language by building one system from the first line to the last, instead of presenting a sequence of unrelated exercises. The system is a laboratory tool called \golab{}. In its first form it is a command-line program that finds files in a research directory, reads them and records their cryptographic digests, runs external programs on them, and writes each of these actions into a ledger whose entries are chained together so that later alteration can be detected. Later parts of the book give the same tool a concurrent worker pool, durable indexes, an HTTP interface, container packaging, and a set of experimental harnesses. At every stage the program that exists is the program the chapter describes, and it compiles, passes its tests, and runs.

The conventional route into Go begins with variables and loops and arrives at web services by way of a long tour of the standard library. The route taken here is different because the intended reader is someone who already works fluently at a shell, who has written Bash scripts that traverse directories and invoke converters, and who has felt the point at which such scripts stop being adequate. The point is rarely a matter of speed. It arrives when a workflow must keep distinctions that the shell leaves implicit: which files were merely found and which were actually read, which outputs are evidence and which are interpretations of evidence, which failures left the filesystem partly changed, and which actions were permitted rather than merely performed. Go is a good language in which to write such distinctions down, because its types are simple enough to be read quickly, its errors are ordinary values that must be handled in plain sight, and its tools produce a single executable that can be run from the same shell the reader already uses.

Several ideas that organize the author's wider research recur here as questions that shape the design of \golab{}: the difference between a proposed transition and an accepted one, the difference between recording an event and authorizing it, the difference between a source and a summary of that source, and the difference between routing work to a destination and adjudicating whether the work should be done. These ideas are used as structural analogies and as prompts for design. The language does not enforce any of them automatically, and the book takes care to say, at each point, whether a given property is guaranteed by the compiler, guaranteed by a library under stated assumptions, or merely maintained by convention and discipline. A reader who finds that a chapter makes a stronger claim than its code supports should treat that as a defect of the chapter.

The book is organized in six parts. Part~I, Local Computation, runs through six chapters. It introduces programs as transition systems, values and types, filesystem traversal and hashing, errors as values, the execution of external programs, and the hash-chained ledger. Part~II, Governed Workflows, runs through three further chapters and introduces interfaces and composition, together with a five-stage record pipeline in which storing, validating, admitting, projecting, and activating a record are separate steps with separate types. Part~III, Concurrent Execution, runs through three more chapters and gives the tool a bounded worker pool, cancellation through contexts, and a ledger writer that serializes appends. Part~IV, Memory and Reconstruction, runs through three more chapters and adds a file that is replaced atomically and an index that is derived from the ledger and can always be rebuilt. Part~V, Distributed Systems, runs through three more chapters and gives the tool an HTTP interface, separates routing from adjudication and from admission, and runs the service in a minimal container. Part~VI, Experimental Systems, runs through four chapters and treats the tool as an object of study, with benchmarks and profiles, structured logs and injected faults, manifests that fix what an experiment ran on, and a search for what a summary of a ledger discards. All six parts are complete in this volume.

Every listing in Parts~I to VI is included directly from the source files of the \golab{} module, so the text and the code cannot drift apart. The module was built and tested with Go 1.24 on Linux, and nothing in the whole book depends on a feature newer than Go 1.21. Readers working under the Windows Subsystem for Linux can follow the book without modification. Commands that appear in a shell listing were run, and the output shown is the output that was produced. Appendix~\ref{app:guarantees} collects, in one table, which properties the compiler or toolchain guarantees and which are conventions, so that no claim in the text need be read as stronger than that table allows.
